%───────────────────────────────────────────────────────────────────────────────
\chapter{Código y algoritmos}
\label{ch:codigo}
%───────────────────────────────────────────────────────────────────────────────

El template distingue dos tipos de material computacional con herramientas
distintas. El paquete \texttt{listings}~\cite{listings2020} es para código
fuente ejecutable real: produce syntax highlighting completo y puede importar
archivos directamente desde el sistema de archivos. El paquete
\texttt{algorithm2e}~\cite{algorithm2e2017} es para pseudocódigo matemático
independiente del lenguaje: opera en modo texto con macros de coloreado
semántico y ecuaciones.

La distinción no es estética sino epistémica: un bloque de listings dice
\textit{``este código corre en Python 3.11''}; un bloque de algorithm2e
dice \textit{``este procedimiento existe en el nivel matemático, su
implementación es ortogonal a la exposición''}.

%───────────────────────────────────────────────────────────────────────────────
\section{Listings: código fuente ejecutable}
\label{sec:listings}
%───────────────────────────────────────────────────────────────────────────────

\subsection{Sintaxis básica}
\label{subsec:listings-sintaxis}

El bloque de código más simple usa un estilo predefinido y un caption:

\begin{lstlisting}[style=latex, frame=none, numbers=none]
\begin{lstlisting}[style=python, caption={Nombre descriptivo}, label={lst:etiqueta}]
# codigo aqui
end{lstlisting}   % <- sin backslash en el documento real
\end{lstlisting}

Los estilos disponibles son: \texttt{python}, \texttt{c}, \texttt{cpp},
\texttt{rust}, \texttt{go}, \texttt{julia}, \texttt{matlab},
\texttt{fortran}, \texttt{asm}, \texttt{mathematica}, y \texttt{latex}.
Todos heredan del estilo base \texttt{catppuccin-base} definido en
\texttt{config/listings.tex}.

\subsection{Importar archivos externos}
\label{subsec:lstinputlisting}

\begin{definicion}{lstinputlisting}
\label{def:lstinputlisting}
El comando \texttt{\textbackslash lstinputlisting} importa un archivo de
código directamente desde el sistema de archivos, sin necesidad de copiar
el contenido en el \texttt{.tex}. La carpeta \texttt{code/} del proyecto
es el repositorio canónico de todos los archivos de código.
\end{definicion}

\begin{lstlisting}[style=latex, frame=none, numbers=none]
% Archivo completo
\lstinputlisting[style=python,
  caption={Integrador de Euler para EDOs},
  label={lst:euler}]{code/euler.py}

% Fragmento: solo lineas 15-40
\lstinputlisting[style=python,
  firstline=15, lastline=40,
  caption={Funcion de paso del integrador},
  label={lst:euler-step}]{code/euler.py}

% Con numero de lineas relativo al archivo
\lstinputlisting[style=julia,
  firstnumber=15,
  firstline=15, lastline=40,
  caption={...}]{code/rk4.jl}
\end{lstlisting}

\begin{observacion}{Rutas relativas}
Las rutas en \texttt{\textbackslash lstinputlisting} son relativas al
archivo \texttt{main.tex}, no al archivo \texttt{.tex} del capítulo.
Si \texttt{main.tex} está en la raíz del proyecto, \texttt{code/euler.py}
apunta a \texttt{./code/euler.py} desde la raíz, independientemente de
en qué capítulo aparezca el comando.
\end{observacion}

%───────────────────────────────────────────────────────────────────────────────
\section{Galería de estilos por lenguaje}
\label{sec:galeria}
%───────────────────────────────────────────────────────────────────────────────

Los snippets siguientes son fragmentos reales de física computacional y
machine learning. Cada uno está elegido para ejercitar las distintas clases
de color: keywords estructurales, tipos, funciones, strings, comentarios,
macros y módulos.

\subsection{Python — Descenso de gradiente estocástico}
\label{subsec:python}

\begin{lstlisting}[style=python,
  caption={SGD con momentum para un modelo de regresión logística},
  label={lst:sgd}]
import numpy as np
from typing import Callable

def sgd_momentum(
    grad_fn: Callable[[np.ndarray], np.ndarray],
    theta0:  np.ndarray,
    lr:      float = 1e-3,
    beta:    float = 0.9,
    n_iter:  int   = 1000,
) -> np.ndarray:
    """Descenso de gradiente estocastico con momentum de Polyak."""
    theta = theta0.copy()
    v     = np.zeros_like(theta)   # momento acumulado

    for t in range(1, n_iter + 1):
        g  = grad_fn(theta)        # gradiente estocastico
        v  = beta * v + (1 - beta) * g
        theta -= lr * v            # actualizacion de parametros

    return theta
\end{lstlisting}

\subsection{C — Integrador de Euler para una EDO escalar}
\label{subsec:c}

\begin{lstlisting}[style=c,
  caption={Método de Euler explícito: $\equ{y}' = \eqfn{f}(\eqk{t}, \equ{y})$},
  label={lst:euler-c}]
#include <stdio.h>
#include <math.h>

/* f(t, y) = -lambda * y  =>  solucion exacta: y(t) = y0 * exp(-lambda*t) */
static double f(double t, double y, double lam) {
    (void)t;
    return -lam * y;
}

void euler(double y0, double t0, double tf,
           double lam, int n, double *out_t, double *out_y) {
    double h = (tf - t0) / n;
    out_t[0] = t0;
    out_y[0] = y0;

    for (int i = 0; i < n; ++i) {
        out_y[i+1] = out_y[i] + h * f(out_t[i], out_y[i], lam);
        out_t[i+1] = out_t[i] + h;
    }
}
\end{lstlisting}

\subsection{C++ — Resolución de sistema lineal con LU (Eigen)}
\label{subsec:cpp}

\begin{lstlisting}[style=cpp,
  caption={Factorización LU parcialmente pivotada con Eigen},
  label={lst:lu-cpp}]
#include <Eigen/Dense>
#include <iostream>

using Matrix = Eigen::MatrixXd;
using Vector = Eigen::VectorXd;

// Resuelve A x = b mediante descomposicion LU con pivoteo parcial.
// Lanza std::runtime_error si A es singular numericamente.
Vector solve_lu(const Matrix& A, const Vector& b) {
    auto lu = A.partialPivLu();
    if (std::abs(lu.determinant()) < 1e-12)
        throw std::runtime_error("Matriz singular: det(A) ~ 0");
    return lu.solve(b);
}

int main() {
    Matrix A = Matrix::Random(4, 4);
    Vector b = Vector::Random(4);
    Vector x = solve_lu(A, b);
    std::cout << "Residuo: " << (A * x - b).norm() << "\n";
}
\end{lstlisting}

\subsection{Rust — Integrador RK4 genérico}
\label{subsec:rust}

\begin{lstlisting}[style=rust,
  caption={Runge-Kutta de orden 4 genérico sobre \texttt{f64}},
  label={lst:rk4-rust}]
/// Avance de un paso del metodo RK4.
/// `f`: derivada del sistema  dy/dt = f(t, y)
/// Devuelve y(t + h).
pub fn rk4_step<F>(f: F, t: f64, y: f64, h: f64) -> f64
where
    F: Fn(f64, f64) -> f64,
{
    let k1 = f(t,           y);
    let k2 = f(t + 0.5*h,  y + 0.5*h*k1);
    let k3 = f(t + 0.5*h,  y + 0.5*h*k2);
    let k4 = f(t + h,       y + h*k3);
    y + (h / 6.0) * (k1 + 2.0*k2 + 2.0*k3 + k4)
}

#[cfg(test)]
mod tests {
    use super::*;
    #[test]
    fn decay_exact() {
        // y' = -y, y(0) = 1  =>  y(1) = e^{-1}
        let f = |_t: f64, y: f64| -y;
        let y1 = rk4_step(f, 0.0, 1.0, 1.0);
        assert!((y1 - (-1.0_f64).exp()).abs() < 1e-6);
    }
}
\end{lstlisting}

\subsection{Go — Caminata aleatoria 2D}
\label{subsec:go}

\begin{lstlisting}[style=go,
  caption={Simulación de caminata aleatoria 2D en Go},
  label={lst:rw-go}]
package main

import (
    "fmt"
    "math/rand"
)

type Point struct{ X, Y float64 }

// RandomWalk simula n pasos de una caminata aleatoria 2D.
// Cada paso tiene longitud 1 en una direccion aleatoria.
func RandomWalk(n int, seed int64) []Point {
    rng   := rand.New(rand.NewSource(seed))
    path  := make([]Point, n+1)
    path[0] = Point{0, 0}
    steps := [][2]float64{{1,0},{-1,0},{0,1},{0,-1}}

    for i := range n {
        s       := steps[rng.Intn(4)]
        path[i+1] = Point{path[i].X + s[0], path[i].Y + s[1]}
    }
    return path
}

func main() {
    path := RandomWalk(10_000, 42)
    fmt.Printf("Posicion final: (%.2f, %.2f)\n",
        path[len(path)-1].X, path[len(path)-1].Y)
}
\end{lstlisting}

\subsection{Julia — Método de potencias para el eigenvalor dominante}
\label{subsec:julia}

\begin{lstlisting}[style=julia,
  caption={Iteración de potencias: eigenvalor dominante de $\eqten{A}$},
  label={lst:power-julia}]
using LinearAlgebra

"""
    power_iteration(A, tol=1e-10, maxiter=1000) -> (lambda, v)

Calcula el eigenvalor dominante lambda y su eigenvector normalizado v
de la matriz simetrica A mediante iteracion de potencias.
"""
function power_iteration(A::AbstractMatrix{T};
                         tol::Float64 = 1e-10,
                         maxiter::Int = 1000) where T <: AbstractFloat
    n = size(A, 1)
    v = normalize(randn(n))      # vector inicial aleatorio
    lambda_old = zero(T)

    for _ in 1:maxiter
        w      = A * v           # multiplicacion matriz-vector
        lambda = dot(v, w)       # cociente de Rayleigh
        v      = normalize(w)
        abs(lambda - lambda_old) < tol && break
        lambda_old = lambda
    end
    return lambda_old, v
end
\end{lstlisting}

\subsection{MATLAB — Ecuación de calor 1D por diferencias finitas}
\label{subsec:matlab}

\begin{lstlisting}[style=matlab,
  caption={Ecuación de calor 1D: esquema de Crank-Nicolson [LeVeque 2007]},
  label={lst:heat-matlab}]
function U = heat_crank_nicolson(alpha, L, T, Nx, Nt)
% heat_crank_nicolson  Resuelve u_t = alpha * u_xx en [0,L] x [0,T]
%   Condiciones de frontera: u(0,t) = u(L,t) = 0
%   Condicion inicial:       u(x,0) = sin(pi*x/L)

dx = L / Nx;  dt = T / Nt;
r  = alpha * dt / (2 * dx^2);    % numero de Courant-Friedrich-Lewy

% Matrices tridiagonales del esquema implicito
e = ones(Nx-1, 1);
A = spdiags([-r*e, (1+2*r)*e, -r*e], [-1 0 1], Nx-1, Nx-1);
B = spdiags([ r*e, (1-2*r)*e,  r*e], [-1 0 1], Nx-1, Nx-1);

x = linspace(dx, L-dx, Nx-1)';
U = sin(pi * x / L);             % condicion inicial

for n = 1:Nt
    U = A \ (B * U);             % paso de tiempo Crank-Nicolson
end
end
\end{lstlisting}

\subsection{Fortran — Multiplicación matriz-vector con BLAS}
\label{subsec:fortran}

\begin{lstlisting}[style=fortran,
  caption={Llamada a \texttt{dgemv} de BLAS para $\eqvec{\mathbf{y}} \leftarrow \alpha\eqten{A}\eqvec{\mathbf{x}} + \beta\eqvec{\mathbf{y}}$},
  label={lst:blas-fortran}]
program matvec_blas
  implicit none
  integer,  parameter :: n = 4
  real(8),  parameter :: alpha = 1.0d0, beta = 0.0d0
  real(8)             :: A(n,n), x(n), y(n)
  integer             :: i

  ! Inicializar A y x
  A = reshape([(real(i,8), i=1,n*n)], [n,n])
  x = [(1.0d0, i=1,n)]
  y = 0.0d0

  ! y <- alpha*A*x + beta*y  (BLAS nivel 2)
  call dgemv('N', n, n, alpha, A, n, x, 1, beta, y, 1)

  write(*,'(A,*(F8.2))') 'y = ', y
end program matvec_blas
\end{lstlisting}

\subsection{Assembler x86-64 — Norma euclídea con SIMD}
\label{subsec:asm}

\begin{lstlisting}[style=asm,
  caption={Norma euclídea de un vector float con instrucciones SSE4},
  label={lst:norm-asm}]
; norm_f32(float* v, int n) -> float
; Calcula ||v||_2 = sqrt(sum v[i]^2) usando SSE4.1
; Argumento: rdi = puntero, esi = longitud (multiplo de 4)
section .text
global norm_f32
norm_f32:
    xorps   xmm0, xmm0          ; acumulador = 0.0
.loop:
    movups  xmm1, [rdi]         ; cargar 4 floats
    mulps   xmm1, xmm1          ; xmm1 = v[i]^2
    addps   xmm0, xmm1          ; acumular
    add     rdi, 16             ; avanzar 4 * sizeof(float)
    sub     esi, 4
    jg      .loop               ; continuar si quedan elementos

    haddps  xmm0, xmm0          ; reduccion horizontal
    haddps  xmm0, xmm0
    sqrtss  xmm0, xmm0          ; sqrt del escalar
    ret
\end{lstlisting}

\subsection{Mathematica — Solución simbólica de la ecuación de onda}
\label{subsec:mathematica}

\begin{lstlisting}[style=mathematica,
  caption={Solución simbólica de la ecuación de onda 1D con condiciones de Dirichlet},
  label={lst:wave-mathematica}]
(* Ecuacion de onda: u_tt = c^2 u_xx en [0,L], u(0,t)=u(L,t)=0 *)
With[{c = 1, L = Pi},
  u[x_, t_] = Sum[
    With[{omega = n*c*Pi/L},
      (A[n]*Cos[omega*t] + B[n]*Sin[omega*t]) * Sin[n*Pi*x/L]],
    {n, 1, 10}];

  (* Condicion inicial: u(x,0) = sin(x), u_t(x,0) = 0 *)
  coefs = Table[
    {n, 2/L * Integrate[Sin[x]*Sin[n*Pi*x/L], {x,0,L}]},
    {n, 1, 10}];

  Plot[u[x, 0.5] /. {A[n_] :> coefs[[n,2]], B[n_] :> 0},
       {x, 0, L}, PlotLabel -> "u(x, 0.5)"]
]
\end{lstlisting}

%───────────────────────────────────────────────────────────────────────────────
\section{Algoritmos: pseudocódigo matemático}
\label{sec:algoritmos}
%───────────────────────────────────────────────────────────────────────────────

\begin{definicion}{Cuándo usar algorithm2e vs listings}
\label{def:algo-vs-listings}
\texttt{algorithm2e} es el entorno correcto cuando el procedimiento existe
en el nivel matemático — cuando su descripción no depende de ningún lenguaje
de programación específico y cuando los símbolos del pseudocódigo son objetos
matemáticos con sus propios roles semánticos. Si el código se va a ejecutar
tal cual, usar \texttt{listings}.
\end{definicion}

\subsection{Sintaxis básica de algorithm2e}
\label{subsec:algo-sintaxis}

\begin{lstlisting}[style=latex, frame=none, numbers=none]
\begin{algorithm}[H]
  \DontPrintSemicolon
  \SetKwInOut{Input}{Entrada}\SetKwInOut{Output}{Salida}
  \Input{descripcion de entradas}
  \Output{descripcion de salidas}

  inicializacion\;
  \While{condicion}{
    cuerpo del bucle\;
    \If{condicion}{
      rama verdadera\;
    }
  }
  \Return{resultado}
  \caption{Nombre del algoritmo [clave-bib]}
  \label{alg:etiqueta}
\end{algorithm}
\end{lstlisting}

Los estilos de texto \texttt{\textbackslash algkw\{\}}, \texttt{\textbackslash algarg\{\}},
\texttt{\textbackslash algcomment\{\}} definidos en \texttt{config/algorithms.tex}
se aplican automáticamente vía \texttt{\textbackslash SetKwSty} y compañía.
No es necesario usarlos manualmente en el cuerpo del algoritmo.

\subsection{Theming semántico de ecuaciones en algoritmos}
\label{subsec:algo-ecuaciones}

La intersección más poderosa del template: pseudocódigo matemático donde
cada símbolo lleva su rol semántico. Las macros de \texttt{equations.tex}
funcionan dentro de \texttt{algorithm2e} en modo matemático normal.

\begin{algorithm}[H]
  \DontPrintSemicolon
  \SetKwInOut{Input}{Entrada}\SetKwInOut{Output}{Salida}
  \Input{%
    NFA $\eqk{N} = (\eqdom{Q}, \eqdom{\Sigma}, \eqop{\delta},
    \eqk{q_0}, \eqdom{F})$}
  \Output{%
    DFA $\eqk{D} = (\eqdom{Q'}, \eqdom{\Sigma}, \eqop{\delta'},
    \eqk{q'_0}, \eqdom{F'})$}

  $\eqk{q'_0} \gets \eqop{E}(\eqdom{\{}\eqk{q_0}\eqdom{\}})$\;
  Añadir $\eqk{q'_0}$ a $\eqdom{Q'}$ como estado no marcado\;

  \While{$\exists\, \equ{T} \in \eqdom{Q'}$ no marcado}{
    Marcar $\equ{T}$\;
    \For{cada $\equ{a} \in \eqdom{\Sigma}$}{
      $\equ{U} \gets \eqop{E}(\eqop{move}(\equ{T}, \equ{a}))$\;
      \If{$\equ{U} \notin \eqdom{Q'}$}{
        Añadir $\equ{U}$ a $\eqdom{Q'}$ no marcado\;
      }
      $\eqop{\delta'}(\equ{T}, \equ{a}) \gets \equ{U}$\;
    }
  }
  \Return{$\eqk{D}$}
  \caption{Construcción de subconjuntos: NFA $\to$ DFA~\cite{press2007numerical}}
  \label{alg:subset-construction}
\end{algorithm}

El color comunica la estructura lógica del algoritmo: $\eqdom{Q}$,
$\eqdom{\Sigma}$, $\eqdom{F}$ son conjuntos (dominios, sapphire).
$\eqk{q_0}$, $\eqk{N}$, $\eqk{D}$ son parámetros conocidos (mauve).
$\equ{T}$ y $\equ{U}$ son variables de trabajo (incógnitas locales, peach).
$\eqop{\delta}$ y $\eqop{E}$ son operadores/funciones (blue, green).

\subsection{Algoritmo numérico: gradiente conjugado}
\label{subsec:conjugate-gradient}

\begin{algorithm}[H]
  \DontPrintSemicolon
  \SetKwInOut{Input}{Entrada}\SetKwInOut{Output}{Salida}
  \Input{%
    Matriz $\eqk{A} \in \eqdom{\mathbb{R}^{n\times n}}$ SPD,
    vector $\eqk{b} \in \eqdom{\mathbb{R}^n}$,
    tolerancia $\eqk{\epsilon} > 0$}
  \Output{Solución aproximada $\equ{x} \approx \eqk{A}^{-1}\eqk{b}$}

  $\equ{x} \gets \mathbf{0}$\;
  $\equ{r} \gets \eqk{b} - \eqk{A}\equ{x}$\;
  $\equ{p} \gets \equ{r}$\;
  $\equ{\rho} \gets \equ{r}^\top \equ{r}$\;

  \While{$\sqrt{\equ{\rho}} > \eqk{\epsilon}$}{
    $\equ{q} \gets \eqk{A}\equ{p}$\;
    $\equ{\alpha} \gets \equ{\rho}\, /\, (\equ{p}^\top \equ{q})$\;
    $\equ{x} \gets \equ{x} + \equ{\alpha}\,\equ{p}$\;
    $\equ{r} \gets \equ{r} - \equ{\alpha}\,\equ{q}$\;
    $\equ{\rho_{\mathrm{new}}} \gets \equ{r}^\top \equ{r}$\;
    $\equ{p} \gets \equ{r} + (\equ{\rho_{\mathrm{new}}} / \equ{\rho})\,\equ{p}$\;
    $\equ{\rho} \gets \equ{\rho_{\mathrm{new}}}$\;
  }
  \Return{$\equ{x}$}
  \caption{Método del gradiente conjugado~\cite{higham2002accuracy}}
  \label{alg:cg}
\end{algorithm}

\begin{observacion}{Algoritmos largos y ruptura de página}
\texttt{algorithm2e} no envuelve el entorno \texttt{algorithm} en
\texttt{tcolorbox} (ver \cref{ch:estructura}). Esto permite que algoritmos
largos se dividan naturalmente entre páginas. Usar el placement \texttt{[H]}
solo cuando el algoritmo es corto y debe aparecer junto al texto que lo
introduce. Para algoritmos de más de 20 líneas, dejar el placement vacío y
permitir que \LaTeX{} lo posicione.
\end{observacion}

%───────────────────────────────────────────────────────────────────────────────
\section{Referenciando código y algoritmos}
\label{sec:ref-codigo}
%───────────────────────────────────────────────────────────────────────────────

Con \texttt{cleveref} configurado en \texttt{config/packages.tex}:

\begin{lstlisting}[style=latex, frame=none, numbers=none]
% Referenciar un bloque de listing
Como muestra el \cref{lst:sgd}, el momentum...

% Referenciar un algoritmo
El \cref{alg:cg} converge en a lo sumo $n$ iteraciones...

% Rango de listings
Los \crefrange{lst:euler-c}{lst:rk4-rust} ilustran...
\end{lstlisting}

Los nombres en español ya están configurados: \texttt{\textbackslash cref\{lst:sgd\}}
produce \textit{cód.~1} y \texttt{\textbackslash Cref\{alg:cg\}} produce
\textit{Algoritmo~2}.

%───────────────────────────────────────────────────────────────────────────────
\printbibliography[heading=subbibliography]
